\section{Literature Review}\label{sec:literature}

This section reviews the literature on extreme-value theory for environmental data, spatial extremes and max-stable processes, pairwise extremal-dependence measures, data-generating processes for the simulation study, inference and uncertainty for dependence estimators, and seasonal and non-stationary extremal dependence.

\subsection{Extreme-value foundations for environmental data}\label{sec:lit-evt}

Extreme-value theory provides the statistical framework for studying rare events in the tails of a distribution. For the canonical case of maxima, let $X_1,\ldots,X_n$ be independent observations with common distribution function $F$ and let $M_n=\max\{X_1,\ldots,X_n\}$. Then
\[
    \Pr\{M_n \leq z\}
    = \Pr\{X_1 \leq z,\ldots,X_n \leq z\}
    = F(z)^n .
\]
The results of \citet{fisher1928} and \citet{gnedenko1943} show that, under normalisation, any non-degenerate limit of normalised sample maxima must belong to the generalised extreme value (GEV) family. The Gumbel, Fr\'echet, and Weibull types are special cases of the GEV distribution, whose shape parameter governs tail behaviour. This result justifies modelling tail observations directly rather than average behaviour, and it provides the theoretical basis for the block-maxima approach, in which the data are divided into non-overlapping blocks and the maximum within each block is modelled. The block-maxima method is also the entry point for spatial extremes, because spatial max-stable processes arise as the limiting objects for spatial block maxima \citep{dehaan2006}.

In our paper the blocks are days: the hourly wind-gust observations are aggregated to daily maximum wind gusts at each station, and these daily maximum wind gusts are the data object on which all subsequent dependence estimation is performed. Daily maximum wind gusts are a deliberate compromise on the block length. On the one hand, raw hourly gusts are strongly serially dependent, since consecutive hours typically belong to the same passing storm system; aggregating to a daily maximum removes the most immediate within-day temporal clustering and yields blocks that are much closer to the independent-maxima setting that the GEV theory describes. On the other hand, taking the block to be a whole season would leave only one observation per station per season, that is roughly $35$ seasonal maxima over $1991$--$2025$, which is far too few to estimate pairwise spatial dependence stably for each of the many station pairs. Daily maximum wind gusts retain on the order of ninety observations per station per season, a sample large enough for the pairwise estimators below while still respecting the block-maxima logic.

A second method is the peaks-over-threshold (POT) framework, which studies observations exceeding a high threshold rather than block maxima. Its theoretical justification is the Pickands--Balkema--de Haan theorem: exceedances above a sufficiently high threshold can be approximated by a generalised Pareto distribution \citep{balkema1974,pickands1975}. POT uses more tail information than block maxima \citep{brabson2000}, but its threshold is finite in observed data while the theory concerns an asymptotic approximation, so the threshold choice is itself an inferential decision \citep{coles2001}. A further difficulty for hourly data is that consecutive exceedances usually belong to the same storm and cannot be treated as independent tail events; correcting for this requires declustering, for instance the automatic interexceedance-time scheme of \citet{ferro2003}, or explicit temporal modelling of the exceedances, as in the Markov-chain treatment of hourly wind speeds by \citet{fawcett2006}. Because daily aggregation already addresses the most immediate form of this serial dependence, our paper treats POT and declustering as background and as a possible robustness alternative rather than as the main empirical strategy.

The marginal extreme-value step is only the first layer of the research cake. The second layer is spatial: conditional on one station observing an extreme daily maximum wind gust, how likely is it that another station also observes an extreme value on the same day? This shifts the focus from univariate extremes to spatial extremal dependence, whose estimands are pairwise tail-dependence measures, extremal coefficients, and dependence--distance curves, which we review in the next sections.

\subsection{Spatial extremes and max-stable processes}\label{sec:lit-spatial}

For block maxima observed over space, max-stable processes provide the standard limiting framework. A max-stable process arises as the limit of pointwise maxima of independent replications of a spatial random field, similar to how the GEV distribution arises as the limit of scalar maxima. Within this framework extremal dependence becomes a spatial object: the joint upper tail is described as a function of location and distance. \citet{dehaan2006} give a foundational treatment of stationary max-stable models for spatial extremes, showing how spatial extremal dependence can be governed by a small number of parameters and related to the distance between locations. The limitation is that stationary max-stable models impose a homogeneous dependence structure, which may be restrictive for Dutch gusts, since winter storms and summer storms do not necessarily share the same spatial dependence pattern.

\citet{davison2012} review the main statistical approaches to spatial extremes, including latent-variable models, copula models, and max-stable processes. The conclusion from their comparison is that good marginal modelling is not sufficient for modelling spatial extremes: models that fit univariate extremes well may still misrepresent joint extremes across locations. This motivates our focus on extremal-dependence summaries rather than on marginal gust behaviour alone. Likelihood-based inference for max-stable processes is hard, because the full multivariate density is either unavailable or computationally infeasible once the number of locations is moderate. \citet{padoan2010} replace the full likelihood with a pairwise composite likelihood, making parametric max-stable inference feasible for spatial extremes and establishing pairwise inference as a standard methodological mitigation.

In our paper, max-stable processes serve mainly as a theoretical reference point and as the benchmark data-generating process for the simulation study (Section~\ref{sec:lit-sim}), rather than as a model that is fully fitted to the Dutch wind-gust data. We adopt the pairwise logic of \citet{padoan2010}, but the empirical objects are nonparametric pairwise summaries computed directly from the seasonal daily maximum wind gusts, not the parameters of a fitted max-stable family. This keeps the winter--summer comparison transparent and avoids committing the empirical analysis to a single parametric dependence family across inter-station distances.

\subsection{Pairwise extremal-dependence measures}\label{sec:lit-pairwise}

The empirical analysis requires dependence measures that are interpretable at the station-pair level. \citet{cooley2012} survey several tools for measuring spatial extremal dependence, including extremal coefficients, madograms, and tail-dependence summaries. These reduce the high-dimensional spatial dependence problem to interpretable pairwise quantities that can be compared across seasons and plotted against inter-station distance. In our paper all three are estimated from the seasonal daily maximum wind gusts, separately for winter (DJF) and summer (JJA), and the margins are transformed separately by station and season before estimation. Standardising margins station-by-station and season-by-season is important: coastal and inland stations have different local gust climates, and winter and summer differ in their marginal gust distributions, so transforming each station--season margin to a common empirical uniform scale ensures that the winter--summer comparison reflects differences in spatial dependence rather than differences in the marginal distributions.

The extremal coefficient is one of the most common summaries. For two locations $s_i$ and $s_j$, the pairwise extremal coefficient $\theta(s_i,s_j)$ lies in $[1,2]$: values close to one correspond to strong extremal dependence, and values close to two to near-independence. In a seasonal comparison, stronger winter dependence corresponds to lower winter values of $\theta(s_i,s_j)$ than summer values at comparable distances.

An empirical route to the extremal coefficient is the $F$-madogram. \citet{cooley2006} study madogram-based dependence summaries for spatial max-stable random fields. For transformed margins, the $F$-madogram between two locations is
\[
    \nu_F(s_i,s_j)
    = \tfrac{1}{2}\,
    \mathbb{E}\!\left[
        \bigl| F_{s_i}\{X(s_i)\} - F_{s_j}\{X(s_j)\} \bigr|
        \right],
\]
where $F_{s_i}$ and $F_{s_j}$ are the marginal distribution functions at the two locations. Under max-stability, the $F$-madogram and the pairwise extremal coefficient are linked through
\[
    \theta(s_i,s_j)
    = \frac{1 + 2\,\nu_F(s_i,s_j)}{1 - 2\,\nu_F(s_i,s_j)} .
\]
The $F$-madogram gives a direct nonparametric estimate of a dependence measure that can be summarised as a dependence--distance curve, without fitting a parametric max-stable model. The upper-tail dependence coefficient $\chi$ provides an additional interpretation. Let $U(s_i)$ and $U(s_j)$ denote marginally standardised observations on a uniform scale. The asymptotic coefficient is
\[
    \chi(s_i,s_j)
    = \lim_{u \uparrow 1} \Pr\{U(s_j) > u \mid U(s_i) > u\},
\]
provided the limit exists, and the corresponding finite-level quantity is
\[
    \chi_u(s_i,s_j)
    = \Pr\{U(s_j) > u \mid U(s_i) > u\}.
\]
Here the standardised observations are the seasonal daily maximum wind gusts on the uniform scale, so $\chi_u$ is computed within each season from daily co-exceedances. The finite-level $\chi_u$ can be read directly off the research question: if one KNMI station records an extreme daily maximum wind gust, how likely is another station to record one on the same day? Larger values of $\chi_u$ indicate stronger finite-level co-exceedance, just as smaller values of $\theta$ indicate stronger extremal dependence. The coefficients $\chi$ and $\bar\chi$ were introduced to distinguish asymptotic dependence from asymptotic independence by \citet{coles1999}, who also discuss their estimation at finite levels. A caveat is that dependence at observed high levels need not coincide with the limiting asymptotic dependence class. \citet{wadsworth2012} show that max-stable models can be too restrictive when data exhibit dependence at observed levels but become independent in the limit, and \citet{huser2022} make the same point about classical max-stable models: they are theoretically natural for asymptotic dependence, but finite environmental data may display weakening dependence at increasingly high levels. This matters for gusts, since winter storms may be spatially extensive while some summer events are localised. Our paper therefore reports finite-level dependence on a small grid of thresholds rather than selecting a single threshold.

\subsection{Data-generating processes for the simulation study}\label{sec:lit-sim}

To validate whether the estimators can recover a true dependence structure, recovery must be measured against a known dependence--distance relation. The pairwise extremal coefficients of max-stable processes are available in closed form, which makes these processes suitable as benchmark data-generating processes: a simulated field carries a known $\theta(s_i,s_j)$ at every distance, against which the estimated curve can be compared. The composite-likelihood machinery of \citet{padoan2010} is built on exactly these parametric families, and the same families serve here as the known dependence structure to be recovered rather than as an estimation target. The Schlather model \citep{schlather2002} provides a flexible stationary max-stable family with a tractable bivariate dependence function, while the Brown--Resnick process \citep{brown1983}, characterised through a negative-definite variogram, is a widely used family whose extremal dependence is controlled by that variogram and ranges from strong to near-independent as the variogram range is varied. \citet{dombry2016} provide exact simulation algorithms for max-stable processes, which avoid the approximation error of naive truncated constructions and ensure that the known dependence structure is the one actually simulated.

Our paper uses the Brown--Resnick process as the main benchmark data-generating process, because its variogram gives direct and interpretable control over how extremal dependence decays with distance, which is exactly the feature the estimators are meant to recover. The null and alternative designs are implemented by choosing one such family and setting its variogram-range parameters so that the two regimes either share a dependence structure (null) or differ by a controlled amount (alternative), with the alternative calibrated so that one regime is more strongly dependent over distance than the other. We stress that the Brown--Resnick process is used only as a controlled benchmark whose true dependence is known; we do not assume that the Dutch wind-gust data are exactly Brown--Resnick. The purpose of the simulation is solely to test whether the pairwise estimators can recover a known winter-versus-summer difference in spatial extremal dependence under a data-generating process where that difference is fixed by design.

\subsection{Inference and uncertainty for dependence estimators}\label{sec:lit-inference}

Because our paper compares estimated dependence--distance curves, the sampling behaviour of the estimators is itself a concern: a difference between two curves is only informative once the uncertainty in each is quantified. Two sources of variability matter. First, the $F$-madogram and $\chi_u$ estimators are empirical means and conditional frequencies computed from a finite number of seasonal daily maximum wind gusts, so their variance grows as the season-split sample shrinks. Second, daily maximum wind gusts are not fully independent in time: although daily aggregation removes the within-day clustering present in hourly data, a single multi-day storm can produce extreme daily maxima on several consecutive days, so consecutive daily observations may still be dependent. Uncertainty quantification should therefore not treat all daily observations as independent.

For this reason our paper quantifies uncertainty by resampling temporal blocks rather than individual days. A year or seasonal-block bootstrap resamples whole years, or whole seasonal segments, with replacement and recomputes the full dependence--distance curves on each resample, so that the resulting variability reflects both the finite sample size and the residual serial dependence between consecutive daily maximum wind gusts within a storm. This block-resampling logic is the daily-maxima counterpart of the cluster-based resampling used for hourly exceedances: the interexceedance-time and declustering machinery of \citet{ferro2003}, and the temporal modelling of \citet{fawcett2006}, address the same serial-dependence problem at the hourly-exceedance level and remain available as robustness checks. For the spatial estimators, the pairwise composite-likelihood framework of \citet{padoan2010} comes with asymptotic standard errors, but these rely on the parametric model being correctly specified; the nonparametric pairwise summaries used here are instead assessed by block resampling and, in the simulation, by direct Monte Carlo evaluation against the known truth. The simulation study thus plays the role that an analytic sampling distribution would play if one were available, characterising bias and variability as a function of sample size, distance, and threshold.

\subsection{Seasonal and non-stationary extremal dependence}\label{sec:lit-nonstationarity}

The third layer of the research cake can be viewed as a restricted form of non-stationary spatial extremal dependence: instead of one dependence structure for all daily maximum wind gusts, our paper allows the dependence relation to differ between two regimes, with season treated as a simple binary covariate for the dependence structure. \citet{huser2016} develop non-stationary dependence structures for spatial extremes, showing that extremal dependence itself, not only the marginal distribution, can vary with covariates. The seasonal split is a deliberately simple version of this idea, and it also addresses a stationarity concern. Many extreme-value methods rely on stationarity, whereas pooled daily maximum wind gusts combine different meteorological regimes; estimating one dependence curve from the unstratified daily maxima would mix large-scale winter storms with more localised summer convective events. Stratifying by season does not remove all non-stationarity, but it makes the working assumption more credible and gives each estimated dependence--distance curve a clearer interpretation.

The literature above serves distinct methodological roles. The classical results of \citet{fisher1928}, \citet{gnedenko1943}, \citet{balkema1974}, and \citet{pickands1975} justify modelling tail observations and motivate the daily block-maxima choice, while \citet{ferro2003} and \citet{fawcett2006} supply the temporal-dependence corrections that would be needed if the analysis worked from hourly exceedances rather than daily maximum wind gusts. The spatial-extremes work of \citet{dehaan2006}, \citet{davison2012}, and \citet{padoan2010} provides the max-stable and pairwise-inference framework, and \citet{cooley2006} and \citet{cooley2012} motivate the pairwise estimands -- the extremal coefficient, the $F$-madogram, and dependence--distance curves -- estimated here from seasonal daily maximum wind gusts. \citet{coles1999} underpins the finite-level $\chi_u$ diagnostic, while \citet{wadsworth2012} and \citet{huser2022} caution that finite-level dependence may differ from asymptotic dependence, motivating the threshold-grid robustness check. The benchmark data-generating processes for the validation study draw on the max-stable families of \citet{padoan2010} and, in particular, the Brown--Resnick process \citep{brown1983} with exact simulation \citep{dombry2016}.

The applied studies of \citet{bonazzi2012}, \citet{dellamarta2009}, \citet{roberts2014}, \citet{klawa2003}, and \citet{coles1994} supply the context: windstorms are spatial events, gust fields require local standardisation, and wind extremes may depend on storm structure and direction. \citet{huser2016} shows that extremal dependence can itself be non-stationary. The contribution of our paper is to combine these ideas in a deliberately narrow, methodologically focused setting: to establish, by simulation, when a seasonal difference in spatial extremal dependence can be recovered from finite daily maximum wind gusts on a fixed, irregular station network, and then to apply the validated estimators to Dutch wind gusts. Our paper does not build a compound-hazard model, fit a high-dimensional vine copula, or attribute changes to climate change; its object is the finite-sample comparison of winter and summer dependence--distance curves and finite-level tail-dependence probabilities across KNMI station pairs.